\chapter{Appendix}
\label{chap:appendix}

% ─────────────────────────────────────────────────────────────────────────────
\section{CBOM JSON Schema}
\label{sec:cbom_schema}

The following describes the structure of the CBOM JSON file produced by the scanner.
Fields marked \texttt{[opt]} are optional (present only when applicable).

\begin{lstlisting}[language={},caption={CBOM JSON Schema},label=lst:cbom_schema]
{
  "scan_start":  "<ISO-8601 timestamp>",
  "scan_end":    "<ISO-8601 timestamp>",
  "host_count":  <integer>,
  "assets": [
    {
      "host":    "<hostname or IP>",
      "port":    <integer>,

      "connection": {
        "tls_version":    "<TLS 1.0 | 1.1 | 1.2 | 1.3 | unknown>",
        "cipher_suite":   "<cipher name>",
        "verified":       <true | false>,
        "error":          "<error message>"  [opt]
      },

      "key_exchange": {
        "type":      "<pqc_hybrid | pqc_pure | classical | unknown>",
        "algorithm": "<X25519 | X25519Kyber768 | ML-KEM-768 | ...>"
      },

      "certificate": {
        "subject":         "<Distinguished Name>",
        "issuer":          "<Distinguished Name>",
        "algorithm":       "<RSA | ECDSA | ML-DSA | SLH-DSA | ...>",
        "key_bits":        <integer>,
        "type":            "<pqc | rsa | ecdsa | unknown>",
        "not_before":      "<date string>",
        "not_after":       "<date string>",
        "days_until_expiry": <integer>,
        "chain_depth":     <integer>
      },

      "symmetric": {
        "cipher": "<AES-256-GCM | AES-128-GCM | ChaCha20 | 3DES | ...>",
        "class":  "<safe | marginal | broken | unknown>"
      },

      "quantum": {
        "label":           "<FULLY_QUANTUM_SAFE | PQC_READY |
                             NOT_PQC_READY | CRITICAL>",
        "display":         "<[FQS] Fully Quantum Safe | [PQC] PQC Ready |
                             [NOT] Not PQC Ready | [CRIT] Critical>",
        "vulnerabilities": ["<string>", ...],
        "recommendations": ["<string>", ...]
      },

      "scan_timestamp": "<ISO-8601 timestamp>"
    }
  ]
}
\end{lstlisting}

% ─────────────────────────────────────────────────────────────────────────────
\section{Quantum Classification Decision Table}
\label{sec:classification_table}

The following truth table defines the mapping from component classifications
to quantum readiness labels. Priority is evaluated top-to-bottom; the first
matching row determines the label.

\begin{table}[H]
  \centering
  \caption{Classification Decision Table}
  \label{tab:classification}
  \begin{tabularx}{\textwidth}{C{2.2cm} C{2.2cm} C{2.4cm} L{4.2cm} L{3.6cm}}
    \toprule
    \textbf{Key Exchange} & \textbf{Certificate} & \textbf{Symmetric} &
    \textbf{Label} & \textbf{Display} \\
    \midrule
    Any & Any & \texttt{broken} &
      {\footnotesize\texttt{CRITICAL}} &
      \textcolor{critical}{\textbf{\critsym{} Critical}} \\
    \addlinespace
    \texttt{pqc\_*} & \texttt{pqc} & \texttt{safe} &
      {\footnotesize\texttt{FULLY\_QUANTUM\_SAFE}} &
      \textcolor{safe}{\textbf{\fqssym{} FQ Safe}} \\
    \addlinespace
    \texttt{pqc\_*} & \texttt{rsa}/\texttt{ecdsa} & \texttt{safe} &
      {\footnotesize\texttt{PQC\_READY}} &
      \textcolor{pqcready}{\textbf{\pqcsym{} PQC Ready}} \\
    \addlinespace
    \texttt{pqc\_*} & \texttt{pqc} & \texttt{marginal} &
      {\footnotesize\texttt{PQC\_READY}} &
      \textcolor{pqcready}{\textbf{\pqcsym{} PQC Ready}} \\
    \addlinespace
    \texttt{classical} & Any & \texttt{safe}/\texttt{marginal} &
      {\footnotesize\texttt{NOT\_PQC\_READY}} &
      \textcolor{notpqc}{\textbf{\notpqcsym{} Not Ready}} \\
    \addlinespace
    Any & \texttt{rsa}/\texttt{ecdsa} & \texttt{marginal} &
      {\footnotesize\texttt{NOT\_PQC\_READY}} &
      \textcolor{notpqc}{\textbf{\notpqcsym{} Not Ready}} \\
    \addlinespace
    \texttt{unknown} & Any & Any &
      {\footnotesize\texttt{NOT\_PQC\_READY}} &
      \textcolor{notpqc}{\textbf{\notpqcsym{} Not Ready}} \\
    \bottomrule
  \end{tabularx}
\end{table}

% ─────────────────────────────────────────────────────────────────────────────
\section{Environment Variables Reference}
\label{sec:env_vars}

\begin{table}[H]
  \centering
  \caption{Supported Environment Variables}
  \label{tab:env_vars}
  \begin{tabularx}{\textwidth}{L{3.5cm} L{2.5cm} L{8cm}}
    \toprule
    \textbf{Variable} & \textbf{Default} & \textbf{Description} \\
    \midrule
    \texttt{PORT}           & \texttt{5000}  & TCP port on which the Flask
                                               application listens. \\
    \texttt{REPORT\_TTL}    & \texttt{3600}  & Seconds before an in-memory report
                                               is evicted (1 hour). \\
    \texttt{PQC\_CBOM\_OUT} & auto-generated & Full path for the CBOM JSON output
                                               file; set automatically by
                                               \texttt{app.py} when launching
                                               \texttt{scan.sh}. \\
    \bottomrule
  \end{tabularx}
\end{table}

% ─────────────────────────────────────────────────────────────────────────────
\section{Document Revision History}
\label{sec:revision_history}

\begin{table}[H]
  \centering
  \caption{Revision History}
  \label{tab:revision_history}
  \begin{tabularx}{\textwidth}{C{1.5cm} C{2.5cm} L{3cm} L{6.5cm}}
    \toprule
    \textbf{Version} & \textbf{Date} & \textbf{Author} & \textbf{Description} \\
    \midrule
    0.1  & 2026-03-01 & quantumCheck Team & Initial draft — introduction and
                                           scope. \\
    0.5  & 2026-03-10 & quantumCheck Team & Added functional requirements,
                                           architecture, and interface sections. \\
    1.0  & 2026-03-20 & quantumCheck Team & Final draft for hackathon submission;
                                           all sections complete. \\
    \bottomrule
  \end{tabularx}
\end{table}

% ─────────────────────────────────────────────────────────────────────────────
\section{Glossary of Quantum Computing Threats}
\label{sec:quantum_threats}

\begin{description}[noitemsep, leftmargin=2cm]
  \item[Shor's Algorithm] A quantum algorithm (1994) that can factor large
    integers and compute discrete logarithms in polynomial time, breaking
    RSA, ECDSA, ECDH, and Diffie-Hellman. Practical execution requires a
    cryptographically-relevant quantum computer (CRQC) with millions of
    logical qubits.
  \item[Grover's Algorithm] A quantum search algorithm that provides a
    quadratic speedup for unstructured search, effectively halving the
    bit-security of symmetric ciphers. AES-128 is reduced to ~64-bit
    security; AES-256 retains ~128-bit security.
  \item[CRQC] Cryptographically Relevant Quantum Computer — a hypothetical
    future device capable of running Shor's algorithm at scale to break
    current public-key cryptography.
  \item[Harvest Now, Decrypt Later (HNDL)] A threat strategy where adversaries
    capture and store encrypted ciphertext today with the intention of
    decrypting it once a CRQC becomes available. TLS key exchange is the
    primary target.
  \item[Hybrid Key Exchange] A TLS key exchange that combines a classical
    algorithm (e.g.\ X25519) with a PQC algorithm (e.g.\ Kyber) so that
    security holds as long as \textit{either} algorithm remains unbroken.
    Represented by \texttt{X25519Kyber768} in TLS extensions.
\end{description}
